\section{Introduction}~\label{sec:introduction}

Recent advances in robotic systems have significantly improved autonomy in perception, manipulation, and planning across domains such as agriculture and household environments. Despite these advances, robots still frequently fail in real-world environments due to perception errors~\cite{kaipa2016addressing}, action execution failures~\cite{kaelbling1998planning}, and dynamic environmental changes~\cite{kurniawati2016online}. Consequently, many robotic systems still rely on human operators to continuously monitor execution and intervene when failures occur~\cite{dragan2013legibility,goodrich2008human}. However, requiring humans to continuously monitor robotic systems and assess the reliability of robot decisions is impractical for long-horizon human--robot collaboration because it increases cognitive workload and fatigue~\cite{chen2014human}. Instead of relying on humans to identify failures after they occur, robotic systems should be able to recognize uncertain situations on their own and proactively request assistance only when necessary~\cite{sadigh2017active,leusmann2023understanding}. This shift leaves two coupled questions unresolved: (i) \textit{when} should the robot ask a human, and (ii) \textit{what} should it ask?

To answer these questions, the robot must estimate its uncertainty about the world and determine whether additional human input is likely to improve decision making. We therefore formulate robotic decision making as a POMDP~\cite{kaelbling1998planning} and develop a framework that jointly models symbolic knowledge~\cite{davis1993knowledge} and uncertainty within a unified structure. We adopt a STRIPS-style symbolic representation~\cite{fikes1971strips}, where the knowledge base is defined as a set of explicit facts, while uncertainty is represented as a belief distribution over possible world states corresponding to alternative hypotheses about the environment. As observations are collected, the system updates the belief over these states accordingly. By measuring the entropy of this belief distribution, the system estimates uncertainty and determines \textit{when} to query a human. The system also determines \textit{what} to ask by selecting questions that are expected to maximally reduce uncertainty among competing hypotheses. Since interaction only requires simple feedback about the environment state, the proposed framework reduces user cognitive burden without requiring users to make planning decisions.

\placeholderfig{Overview of the proposed knowledge-based robotic system with uncertainty-aware planning. The system maintains a belief over possible worlds and updates it based on observations. If uncertainty exceeds a threshold, the system selectively queries the user to resolve ambiguity; otherwise, it continues autonomous planning and execution.}{1.9in}{fig:overview}

We further design and integrate a system architecture that enables the proposed framework to operate on a real robotic platform. The architecture maintains a belief over possible world states and updates it using both sensor observations and human feedback. A centralized knowledge base connects the sensing, planning, feedback management, and user interface modules, allowing information to be shared consistently across the system. This design enables the robot to reason with the most up-to-date information and request human assistance only when uncertainty arises.

We evaluate the proposed system in two domains, tomato harvesting and waste sorting, using both simulation and real-robot experiments. Experimental results show that the proposed approach reduces the number of human interactions while maintaining or improving task success rates. Furthermore, user studies demonstrate that the system effectively reduces cognitive workload during human--robot collaboration.

The main contributions of this work are as follows:
\begin{itemize}[leftmargin=*]
    \item We propose a knowledge-based robotic framework that jointly represents symbolic world knowledge and uncertainty over possible world states, enabling decision making under partial observability while preserving interpretability.
    \item We present an uncertainty-aware human feedback mechanism that addresses both \textit{when} and \textit{what} to ask by estimating belief uncertainty and selecting informative queries that distinguish competing hypotheses.
    \item We design and integrate a complete robotic system that combines sensing, planning, user feedback, and knowledge management through a centralized knowledge base.
    \item We validate the proposed framework in both tomato harvesting and waste sorting domains through simulation, real-robot experiments, and user studies, demonstrating reduced human intervention and cognitive workload while maintaining task performance.
\end{itemize}
